% !PRINTED 187
\weberpage{187}
%<<CONT
\weq{1}{I(a') \doteq r^{\lambda} I(a)}
besteht. Den Exponenten $\lambda$ wollen wir das \gesperrt{Gewicht} der
Invariante nennen. Ist dieser Exponent Null, so heisst $I$ eine
\gesperrt{absolute Invariante}.

Man betrachtet meist nur \gesperrt{ganze rationale Invarianten},
d. h. solche Functionen $I$, die von den Coëfficienten $a$ rational
abhängen und überdies ganze homogene Functionen von
ihnen sind. Die simultanen Invarianten sollen in Bezug auf die
Coëfficienten einer jeden der Formen, von denen sie abhängen,
homogen sein.

Ausser den Invarianten giebt es auch Functionen, die neben
den Coëfficienten die Variablen selbst noch enthalten, im Uebrigen
aber dieselben Eigenschaften wie die Invarianten (1), „die Invarianteneigenschaft“ haben; also, wenn $C(x, a)$ eine solche
Function ist

\weq{2}{C(x', a') \doteq r^{\lambda} C(x, a).}

Solche Functionen heissen \gesperrt{Covarianten der Form $F(x)$}
% [sic] Im Druck fehlt hier der Satzpunkt nach F(x).
Die Functionaldeterminante ist eine simultane Covariante eines
Systems von $m$ Functionen und die Hesse'sche Determinante
einer Function von höherem als dem zweiten Grade, ist eine
Covariante einer einzelnen Form. Auch unter den Covarianten
betrachtet man vorzugsweise ganze rationale und homogene
Functionen von $x$.

Ist $m$ die Anzahl der Variablen $x$, $n$ der Grad der Function
$F(x)$, $\nu$ und $\mu$ die Grade von $C(x, a)$ in Bezug auf die $x$ und
$a$, so besteht zwischen den Zahlen $\lambda, \mu, \nu, m, n$ die Relation

\weq{3}{n\mu = m\lambda + \nu,}

die man dadurch beweist, dass man den Grad der rechten und
linken Seite von (2) in Bezug auf die Substitutionscoëfficienten
vergleicht. Die $a'$ sind nämlich in den Substitutionscoëfficienten
vom $n^{\text{ten}}$ Grade, $r$ vom $m^{\text{ten}}$ Grade, $x$ vom ersten Grade, woraus
sich (3) ergiebt. Die entsprechende Formel für die Invarianten
erhält man, wenn man $\nu = 0$ setzt, wie denn überhaupt die
Invarianten als specieller Fall der Covarianten betrachtet werden
können.

Ein allgemeines Gesetz zur Bildung von Invarianten und
Covarianten, das wir in den besonderen Fällen des \S.\,59 schon
gebraucht haben, wollen wir kurz besprechen.
